\documentclass[aspectratio=169]{beamer}

\usepackage[utf8]{inputenc}
\usepackage[T1]{fontenc}
\usepackage[brazil]{babel}
\usepackage{amsmath, amssymb}
\usepackage{csquotes}
\usepackage{tikz}

% Colors -- paleta PCC104 (azul escuro + laranja)
\definecolor{mblack}{RGB}{20, 20, 20}
\definecolor{mprimary}{RGB}{30, 45, 47}     % azul escuro
\definecolor{maccent}{RGB}{197, 123, 87}    % laranja
\definecolor{mwhite}{RGB}{255, 255, 255}

% Theme
\usetheme{default}
\usecolortheme{default}
\setbeamercolor{normal text}{fg=mblack, bg=white}
\setbeamercolor{frametitle}{fg=white, bg=mprimary}
\setbeamercolor{title}{fg=white}
\setbeamercolor{author}{fg=white}
\setbeamercolor{date}{fg=white}
\setbeamercolor{itemize item}{fg=maccent}
\setbeamercolor{itemize subitem}{fg=mprimary}
\setbeamercolor{enumerate item}{fg=maccent}
\setbeamercolor{enumerate subitem}{fg=mprimary}
\setbeamercolor{block title}{fg=white, bg=mprimary}
\setbeamercolor{block body}{fg=mblack, bg=mprimary!8}
\setbeamercolor{alerted text}{fg=maccent}
\setbeamercolor{block title alerted}{fg=white, bg=maccent}
\setbeamercolor{block body alerted}{fg=mblack, bg=maccent!10}

\setbeamerfont{frametitle}{series=\bfseries, size=\large}
\setbeamerfont{title}{series=\bfseries, size=\LARGE}
\setbeamerfont{block title}{series=\bfseries}

\setbeamertemplate{navigation symbols}{}

% Footer
\setbeamertemplate{footline}{%
  \leavevmode%
  \hbox{%
    \begin{beamercolorbox}[wd=0.7\paperwidth, ht=2.25ex, dp=1ex, leftskip=0.3cm]{normal text}%
      \tiny\color{mblack!50} DECOM/UFOP -- Boas-vindas: comportamento e \'etica
    \end{beamercolorbox}%
    \begin{beamercolorbox}[wd=0.3\paperwidth, ht=2.25ex, dp=1ex, rightskip=0.3cm, right]{normal text}%
      \tiny\color{mblack!50}\insertframenumber{} / \inserttotalframenumber
    \end{beamercolorbox}%
  }%
}

% Frametitle with accent rule
\setbeamertemplate{frametitle}{%
  \nointerlineskip%
  \begin{beamercolorbox}[wd=\paperwidth, ht=0.55cm, dp=0.2cm, leftskip=0.5cm]{frametitle}%
    \usebeamerfont{frametitle}\insertframetitle%
  \end{beamercolorbox}%
  \vspace{-0.1cm}%
  \textcolor{maccent}{\rule{\paperwidth}{2pt}}%
}

% Title page: three-band layout
\setbeamercolor{titlebar top}{bg=mwhite}
\setbeamercolor{titlebar main}{bg=mprimary}
\setbeamercolor{titlebar bottom}{bg=maccent}

\setbeamertemplate{title page}{%
  \vspace{0pt}%
  \begin{beamercolorbox}[wd=\paperwidth, ht=0.55cm, dp=0pt]{titlebar top}\end{beamercolorbox}%
  \begin{beamercolorbox}[wd=\paperwidth, ht=3.8cm, dp=0pt, center]{titlebar main}%
    \vspace{0.6cm}
    {\usebeamerfont{title}\color{white}\inserttitle\par}%
    \vspace{0.35cm}{\color{white}\small\insertauthor\par}%
    \vspace{0.15cm}{\color{white}\footnotesize\insertdate\par}%
    \vspace{0.25cm}
  \end{beamercolorbox}%
  \begin{beamercolorbox}[wd=\paperwidth, ht=1.3cm, dp=0pt]{titlebar bottom}\end{beamercolorbox}%
}

\setlength{\itemsep}{4pt}

\newcommand{\impt}[1]{\textcolor{maccent}{\textbf{#1}}}
\newcommand{\ntc}[1]{\textcolor{mprimary}{#1}}

\title{Bem-vindos \`a Universidade}
\author{Comportamento acad\^emico \& responsabilidade \'etica em Computa\c{c}\~ao e IA}
\date{\today}

\begin{document}

\begin{frame}[plain]\titlepage\end{frame}

\begin{frame}{Sum\'ario}
\tableofcontents
\end{frame}

% ==========================================================
\section{Parte I -- Comportamento na universidade}
% ==========================================================

\begin{frame}{Voc\^es n\~ao est\~ao mais no ensino m\'edio}
\begin{itemize}
  \item A universidade exige \impt{autonomia}: ningu\'em vai cobrar voc\^es individualmente
  \pause
  \item O professor \ntc{orienta}; quem estuda \'e voc\^e
  \pause
  \item Mais liberdade $\Rightarrow$ \impt{mais responsabilidade}
  \pause
  \item Resultado depende muito mais do que voc\^e faz \emph{fora} da sala
\end{itemize}
\end{frame}

\begin{frame}{Frequ\^encia e presen\c{c}a}
\begin{itemize}
  \item V\'a \`as aulas -- mesmo as ``chatas''
  \pause
  \item \impt{Chegue no hor\'ario}. Atraso atrapalha quem chegou na hora
  \pause
  \item Celular no silencioso, longe da m\~ao
  \pause
  \item Em aula: pergunte. D\'uvida n\~ao perguntada vira nota baixa depois
\end{itemize}
\end{frame}

\begin{frame}{H\'abitos de estudo}
\begin{block}{Regra pr\'atica}
Para cada hora de aula, planeje \impt{2 horas} de estudo fora dela.
\end{block}
\pause
\begin{itemize}
  \item Estude \ntc{antes} da aula seguinte, n\~ao na v\'espera da prova
  \pause
  \item Resolva exerc\'icios -- ler n\~ao \'e estudar
  \pause
  \item Forme grupos de estudo, mas \impt{produza primeiro sozinho}
\end{itemize}
\end{frame}

\begin{frame}{Comunica\c{c}\~ao com professores}
\begin{itemize}
  \item E-mail com \impt{assunto claro} e identifica\c{c}\~ao (nome, matr\'icula, disciplina)
  \pause
  \item Releia antes de enviar. Texto sem pontua\c{c}\~ao n\~ao \'e informalidade -- \'e desleixo
  \pause
  \item \ntc{N\~ao pe\c{c}a} o que est\'a no plano de ensino
  \pause
  \item Hor\'ario de atendimento existe -- use
\end{itemize}
\end{frame}

\begin{frame}{Convivência e respeito}
\begin{itemize}
  \item A universidade \'e \impt{plural}: origens, cren\c{c}as, identidades diferentes
  \pause
  \item Disc\'ordia intelectual: bem-vinda. Desrespeito: n\~ao
  \pause
  \item Trote, ass\'edio, discrimina\c{c}\~ao: \impt{denuncie}
  \pause
  \item Sa\'ude mental importa -- procure o servi\c{c}o de apoio da UFOP quando precisar
\end{itemize}
\end{frame}

\begin{frame}{Integridade acad\^emica}
\begin{alertblock}{O que conta como fraude}
C\'opia, plágio, ``fazer o trabalho do colega'', \impt{usar IA generativa como se fosse seu}, comprar TCC.
\end{alertblock}
\pause
\begin{itemize}
  \item Plágio destr\'oi reputa\c{c}\~ao -- e a sua reputa\c{c}\~ao come\c{c}a agora
  \pause
  \item \ntc{Cite suas fontes}. Inclusive ferramentas de IA, quando autorizadas
  \pause
  \item Em caso de d\'uvida: pergunte ao professor antes de entregar
\end{itemize}
\end{frame}

\begin{frame}{Uso de IA durante o curso}
\begin{itemize}
  \item IA pode ser \impt{aliada} de aprendizagem -- ou atalho para n\~ao aprender
  \pause
  \item Cada disciplina define o que \'e permitido. \ntc{Pergunte}
  \pause
  \item Se voc\^e n\~ao consegue explicar o c\'odigo/texto, \impt{n\~ao \'e seu}
  \pause
  \item Lembrete: a prova \'e sem IA. Quem n\~ao aprendeu vai aparecer
\end{itemize}
\end{frame}

% ==========================================================
\section{Parte II -- \'Etica, sociedade e ambiente}
% ==========================================================

\begin{frame}{Por que isso importa para voc\^es}
\begin{itemize}
  \item Computa\c{c}\~ao e IA \impt{n\~ao s\~ao neutras}
  \pause
  \item O c\'odigo que voc\^es v\~ao escrever \ntc{decide} sobre vidas: cr\'edito, sa\'ude, emprego, justi\c{c}a
  \pause
  \item Engenheiro civil aprende sobre seguran\c{c}a estrutural. Voc\^es precisam aprender sobre \impt{seguran\c{c}a social}
\end{itemize}
\end{frame}

\begin{frame}{Vi\'es e justi\c{c}a algor\'itmica}
\begin{itemize}
  \item Dados refletem desigualdades do mundo real
  \pause
  \item Modelo treinado em dados enviesados \impt{automatiza} a desigualdade
  \pause
  \item Exemplos: reconhecimento facial com m\'a performance em peles negras; sistemas de cr\'edito; tria-gem de curr\'iculos
  \pause
  \item \ntc{Quem n\~ao est\'a nos dados} tamb\'em \'e uma decis\~ao
\end{itemize}
\end{frame}

\begin{frame}{Privacidade e dados pessoais}
\begin{itemize}
  \item LGPD existe -- e \impt{se aplica} ao c\'odigo de voc\^es
  \pause
  \item Coletar \emph{menos} dados costuma ser melhor que coletar mais
  \pause
  \item Anonimiza\c{c}\~ao \ntc{n\~ao} \'e garantia: dados podem ser reidentificados
  \pause
  \item ``Eu n\~ao tenho nada a esconder'' \'e uma m\'a defesa de uma m\'a pol\'itica
\end{itemize}
\end{frame}

\begin{frame}{Trabalho, automa\c{c}\~ao e desigualdade}
\begin{itemize}
  \item IA est\'a transformando profiss\~oes -- inclusive a de voc\^es
  \pause
  \item Quem ganha e quem perde com a automa\c{c}\~ao? \impt{N\~ao \'e dado da natureza}: \'e escolha
  \pause
  \item Concentra\c{c}\~ao de poder em poucas empresas de tecnologia
  \pause
  \item \ntc{Voc\^es t\^em voz} sobre o que constroem e para quem
\end{itemize}
\end{frame}

\begin{frame}{Desinforma\c{c}\~ao e manipula\c{c}\~ao}
\begin{itemize}
  \item Modelos generativos: texto, imagem, voz, v\'ideo (\impt{deepfakes})
  \pause
  \item Custo de produzir mentira convincente caiu drasticamente
  \pause
  \item Impacto em elei\c{c}\~oes, rela\c{c}\~oes pessoais, golpes financeiros
  \pause
  \item Quem desenvolve essas ferramentas \ntc{compartilha} a responsabilidade pelo uso
\end{itemize}
\end{frame}

\begin{frame}{Custo ambiental da computa\c{c}\~ao}
\begin{itemize}
  \item Treinar e operar modelos grandes consome \impt{muita} energia e \'agua
  \pause
  \item Data centers $\Rightarrow$ emiss\~oes de CO$_2$ e estresse h\'idrico local
  \pause
  \item Hardware: mineração, lixo eletr\^onico, obsolesc\^encia programada
  \pause
  \item \ntc{Efici\^encia} (de modelo e de c\'odigo) \'e quest\~ao ambiental, n\~ao s\'o t\'ecnica
\end{itemize}
\end{frame}

\begin{frame}{Seguran\c{c}a e usos maliciosos}
\begin{itemize}
  \item Mesmo c\'odigo bem-intencionado pode ser \impt{reaproveitado} para o mal
  \pause
  \item Vulnerabilidades em sa\'ude, infraestrutura cr\'itica, sistemas financeiros
  \pause
  \item IA aplicada a vigil\^ancia em massa, armas aut\^onomas
  \pause
  \item Pergunta b\'asica: \ntc{quem se machuca} se isso for usado errado?
\end{itemize}
\end{frame}

\begin{frame}{O profissional que voc\^es ser\~ao}
\begin{block}{Tr\^es perguntas antes de codar}
\begin{enumerate}
  \item \impt{Para quem} isso serve? Quem est\'a de fora?
  \item Como isso pode ser usado \impt{para causar dano}?
  \item Eu conseguiria \impt{justificar} essa decis\~ao publicamente?
\end{enumerate}
\end{block}
\pause
\vspace{0.3cm}
\ntc{Recusar} um projeto antiético tamb\'em \'e parte da profiss\~ao.
\end{frame}

\begin{frame}{Mensagem final}
\begin{itemize}
  \item Voc\^es escolheram uma \'area com \impt{enorme poder} -- e portanto enorme responsabilidade
  \pause
  \item Bons profissionais s\~ao t\'ecnicamente competentes \ntc{e} eticamente atentos
\end{itemize}
\end{frame}

\end{document}